\section{What a released past must keep}\label{sec:classes}

The two-bit example of Figure~\ref{fig:three} has a general form, and it is the cleanest place to see the mechanism of the law without any quantum ingredient. A past $x$ arrives and must be released. Later a question $y$ arrives, and the device must answer it with $g(x,y)$. Two things can cost work here, and it pays to keep them apart. One is ordinary erasure: if the answer discards part of $y$, the device must absorb what it discards, whatever the schedule. The other is order: while $x$ is gone, the device must still be able to answer every $y$, so it must keep a record of what $x$ can still tell apart. The theorem below prices both and separates them.

\paragraph{Classes and fibers.} Fix nonempty finite alphabets $X,Y,Z$ and $g:X\times Y\to Z$. Two pasts are \emph{equivalent}, $x\sim x'$, if $g(x,\cdot)=g(x',\cdot)$: no later question distinguishes them. Let $K(g)$ be the number of classes, $c(x)$ the class of $x$, and write $g(c,y)$ for the common value on a class. Two counts measure what an answer discards:
\[
F(g)=\max_{x\in X,\ z\in Z}\bigl|\{y:g(x,y)=z\}\bigr|,
\qquad
F'(g)=\max_{z\in Z}\bigl|\{(c,y):g(c,y)=z\}\bigr| .
\]
$F$ is the most questions a single past leaves indistinguishable after answering; $F'$ counts pairs of a class and a question that lead to the same answer. Always $F\le F'\le K\,F$. A task has \emph{reversible questions} if $y\mapsto g(x,y)$ is a bijection of $Y$ onto $Z$ for every $x$; then $F=1$ and $F'=K$. The classes are the distinct rows of the table of $g$, so $\lceil\log_2K(g)\rceil$ is the deterministic one-way communication complexity of $g$ from the party that holds the past~\cite{KushilevitzNisan96}; read as a strategy, they are its causal states, the memory states of an optimal causal agent in computational mechanics~\cite{BarnettCrutchfield15,Thompson25}. Thompson et al.\ price such agents for a given input process, with each question kept on the tape and each answer written to a fresh cell~\cite{Thompson25}. There the record can be uncomputed against the question and the answer, and the cost they find is the information the memory stores about the last input and output minus what it holds about the next ones~\cite[Eq.~(3)]{Thompson25}, nonpredictive information of the kind that Still, Sivak, Bell and Crooks identified with the average dissipation of a system driven by a stochastic signal~\cite{Still12}; for the task of Figure~\ref{fig:three} that cost is zero. Here the answer overwrites the question, as in the information ratchets of~\cite{BMC18}, and the price is a worst case over observers; Appendix~\ref{app:module} compares the two interfaces.

\begin{theorem}[Fiber law for a released past]\label{thm:classes}
Let the task $T_g$ be: a carrier arrives holding $x$ in the computational basis, is dephased and released; a second carrier, of dimension $\max\{|Y|,|Z|\}$, then arrives holding $y$ and must be released holding $g(x,y)$; its admissible inputs, together with the observer's references, are supported on $\operatorname{span}\{|y\rangle:y\in Y\}$ tensored with the references, a promise that is part of the task. All carriers have zero Hamiltonian. Against quantum devices and observers with references,
\begin{align*}
K_\beta^{S}(T_g)&=\log_2F'(g)\quad\text{bits per episode in real time},\\
K_\beta^{S}(T_g)&=\log_2F(g)\quad\ \text{bits per episode when $x$ may be held until $y$ arrives.}
\end{align*}
The order price, the difference, lies between $0$ and $\log_2K(g)$. For reversible questions it is $\log_2K(g)$ and the held price is zero. For a constant answer $g\equiv z_0$ it is zero, and both schedules pay $\log_2|Y|$: the Landauer cost of erasing $y$.
\end{theorem}

The proof is in Appendix~\ref{app:classes}; its two halves are the two halves of Theorem~\ref{to:thm:law}. The upper bound is a device that dephases each carrier with a coin drawn from the reservoir, writes the class $c(x)$ into a register before $x$ leaves, and when $y$ arrives rewrites $(c,y)$ reversibly as the answer $g(c,y)$ on the carrier and an index of the pair within the fiber of that answer in the register; the index has at most $F'$ values. The lower bound needs two kinds of observer. Observers who might ask a distinguishing question force the device to know the class of $x$ at the moment $x$ leaves, since one fixed later operation must tell the classes apart. The \emph{steering} observer (it steers the answer; the word has nothing to do with quantum steering), who chooses the answer $z_0$ with the largest fiber and sends pasts and questions uniformly over that fiber, then leaves the device no way to pass that knowledge on, because every answer it returns is $z_0$. An identity for the information the device holds about the observer's choices, visit by visit, turns these two facts into an entropy increase of $\log_2F'$ per episode, up to the explicit loss $\delta_g(\varepsilon)$ for sufficiently small $\varepsilon$, uniform in the number of episodes and in the dimension of the private device. The held bound is the thermodynamic capacity of the classical channel $(x,y)\mapsto(x,g(x,y))$, which is $\log_2F$. Both prices are worst-case Landauer costs of one logical operation applied in the same way to every input, $(x,y)\mapsto(x,g(x,y))$ when the past is held and $(c,y)\mapsto g(c,y)$ in real time, the setting of the general forms of Landauer's principle for deterministic and indeterministic operations~\cite{Turgut09,Maroney09}. Merging of this kind has also been priced for finite automata that read their input once, in real time: K{\i}yak and Say charge such an automaton $\log_2P$ bits whenever it reaches a configuration with $P$ possible predecessors, and compare deterministic automata with quantum automata at zero and at bounded error~\cite{KiyakSay24}. Their automaton has a fixed finite memory and is charged step by step for recognizing a language; here the device may have any size, the price is a rate against the worst observer, and the real-time bound holds against every quantum device at small error. What the real-time bound adds is the continuation constraint of a released past, which forces the record of the class, and the pairing of observers that prices that record robustly.

\paragraph{Release rank is the linear case.} When $X=\mathbb F_2^k$, $Y=Z=\mathbb F_2^m$ and $g(x,y)=y+Fx$, the questions are reversible, the classes are the cosets of $\ker F$, and the order price is $\operatorname{rank}(F)$ bits. With several input carriers in flight, some of which may be held, Proposition~\ref{to:prop:rank} gives the price as the rank of the released part $F_L$ and shows that deciding whether holding $m$ carriers lowers it at all is NP-complete.

\paragraph{A record is not yet a price.} The record of $c(x)$ must exist at the moment $y$ arrives, whatever $y$ turns out to be: for pasts in different classes the device's states at that moment are told apart by one fixed later operation, so they are orthogonal, and nearly so at small error (Lemma~\ref{lem:classes-orth}). It is tempting to conclude that the device pays $\log_2K$ against every observer, including one who never asks a distinguishing question. That conclusion is false, and the reason is instructive. For the XOR task the reversible two-gate circuit
\[
(M,Y)=(a,y)\ \longmapsto\ (y,\ a\oplus y)
\]
answers every question correctly and leaves the register holding $y$ rather than $a$: the record of the past has been passed into the answer, as in a reversible computation that uncomputes its history~\cite{Bennett73}. Against an observer who always sends $y=0$ this device ends each episode with an empty register and pays nothing. What the law charges is the entropy the device cannot pass on, and the lower bound pairs the observers who might ask a distinguishing question, which force the record to exist, with the observer who closes the only exit by choosing questions that make the answer constant. Neither suffices alone: for $g(x,y)=x\wedge y$ a device that serves only the steering observer pays $h_2(1/3)\approx0.918$ bits per episode instead of $\log_23$. The price is a property of the worst observer the device must serve, not of the record it must hold.

\paragraph{Order prices are not additive.} One might hope to price a schedule by a sum of terms, one per release decision. For linear tasks the price is a rank function and is submodular in the released set. For reversible tasks in general it is not: in Appendix~\ref{app:classes}, Proposition~\ref{prop:nonsubmodular} gives a task on three bits with a three-valued answer whose order price violates submodularity by $\log_2(3/2)$ bits per episode, the largest for any task of the form $g(x,y)=y+f(x)\bmod3$ on three bits. What is additive is the law itself, across independent episodes, which is why rates are the right object.
